\documentclass[10pt,a4paper]{amsart}
\usepackage[margin=19mm]{geometry}
\usepackage{amsmath,amssymb,mathtools}
\usepackage[scr=rsfs,cal=euler]{mathalfa}
\usepackage{xfrac}
\usepackage[unicode,hidelinks,bookmarks=false]{hyperref}
\newcommand{\ie}{i.e.\ }
\newcommand{\cf}{cf.\ }
\newcommand{\resp}{respectively\ }
\newcommand{\Subsection}{\subsection}
\providecommand{\nobreakdash}{\nobreak\hbox{-}}
\setlength{\emergencystretch}{2em}
\multlinegap0pt
\usepackage[all]{xy}
\usepackage{amssymb}
\usepackage{xcolor}
\newtheorem{lemma}{Lemma}
\usepackage{cleveref}
\crefname{lemma}{Lemma}{Lemmas}
\newcommand{\Adj}{\operatorname{Adj}}%Adjunction
\newcommand{\B}{\mathcal{B}}
\newcommand{\Oone}{\operatorname{O(1)}}
\newcommand{\id}{\operatorname{id}}%identity
\newcommand{\opp}{\operatorname{op}}%opposite
\title{Relative and lax volutive categories}
\author{Tim L\"uders}
\date{Source: arXiv:2602.15667v1 (CC BY 4.0)}
\begin{document}
\maketitle

\noindent\emph{Source and licence.} Relative and lax volutive categories. Version arXiv:2602.15667v1 (CC BY 4.0), distributed by arXiv under the Creative Commons Attribution 4.0 International licence, http://creativecommons.org/licenses/by/4.0/, as recorded in the arXiv licence metadata for this exact version and shown on the abstract page https://arxiv.org/abs/2602.15667v1. Full licence notice: \texttt{licenses/arxiv-2602.15667-CC-BY-4.0.txt}.\par\smallskip

\noindent\textit{Editorial scope.} Counted unit: the article's lemma adjunction (source0.tex lines 958--961). The article's complete original proof of it is delivered with it (source0.tex lines 962--1033, one \texttt{proof} environment of 72 lines). No mathematics is written, repaired or re-derived: the statement and the proof are the article's own lines, in the article's own notation, and no retained line is substituted or re-flowed.  Retained uncounted as the source-local context the proof needs: lines 920--928; lines 930--939.  Nothing was omitted from any retained span, so the package carries no omission record.  No retained line contains a citation, so the wrapper carries no bibliography and no bibliography entry.  No retained line contains a figure environment, an image-inclusion command or drawing code, so no image is delivered and the package declares no asset dependency.  Editorial formatting only, in addition: an AMS \texttt{amsart} preamble that restates the article's own declarations and macros used by the retained lines, namely the environments \texttt{lemma} on separate counters, together with the standard packages those lines need.  The retained environments print in this document as Lemma 1 in order of appearance, whereas the article prints the same items with the numbering of its own full document.  The complete native source of the article as distributed in the arXiv e-print for this exact version is bundled unchanged as \texttt{sources/194812/source0.tex}.
    \begin{equation}\label{Equation: first coherence condition of h in volutive adjunction}
        \begin{aligned}
        \xymatrix@R=0.5em{
            Z \ar[r]^-{\id \circ h} & Z \eta_c^{-1}  d(Z)  Z \ar[r]^-{\eta_Z \circ \id} & \eta_{d(c)}^{-1} d^2(Z)d(Z)Z 
            \ar[r]^-{\tau_c \circ \id} & d(\eta_c^{-1}) d^2(Z) d(Z) Z \\
            \ar[r]_-{\cong} & d(\eta_c^{-1} d(Z)Z) Z \ar[r]_-{d(h) \circ \id} & Z
        }
        \end{aligned}
    \end{equation}

    \begin{equation}\label{Equation: second condition of h in volutive adjunction}
        \begin{aligned}
        \xymatrix@R=0.5em{
            \eta_c^{-1} d(Z) \ar[r]^-{h \circ \id} & \eta_c^{-1} d(Z) Z \eta_c^{-1} d(Z) \ar[r]^-{\id \circ \eta_Z \circ \id} & 
            \eta_c^{-1} d(Z) \eta_{d(c)}^{-1} d^2(Z) d(Z) \\
            \ar[r]_-{\id \circ \tau_c \circ \id} & 
            \eta_c^{-1} d(Z) d(\eta_c^{-1}) d^2(Z) d(Z) \ar[r]_-{\id \circ d(h)} &\eta_c^{-1} d(Z)
        }
        \end{aligned}
    \end{equation}

\begin{lemma}\label{Lemma: volutive adjunctions and volutive functors from the walking adjunction}
    Let $(\B,d,\eta,\tau)$ be a 2-contravariant $\Oone$-volutive 2-category. Then, volutive adjunctions in $(\B,d,\eta,\tau)$ are the same as 
    2-contravariant $\Oone$-volutive 2-functors $\Adj \to \B$. 
\end{lemma}

\begin{proof}
    Let $(c,Z,h)$ be a volutive adjunction in $(\B,d,\eta,\tau)$. We define a functor $G \colon \Adj \to \B$ by setting $G(a) = c, G(b) = d(c), G(X) = Z, G(Y) = \eta_c^{-1} \circ d(Z), G(f) = h, G(g) = d(h) \bullet (\tau_c \circ \id) \bullet (\eta_Z \circ \id)$. 
    We need to check that the Zorro moves are satisfied. The first Zorro move is satisfied by \Cref{Equation: first coherence condition of h in volutive adjunction}
    and the second Zorro move is satisfied by \Cref{Equation: second condition of h in volutive adjunction}. Next, we need to construct a 2-contravariant 
    $\Oone$-volutive structure $(\omega,\Xi)$ on $G$. We define an invertible 2-transformation $\omega \colon G^{2\opp} \circ F \to d \circ G$ as follows. 
    Its 1-morphism component at $a \in \Adj$ is given by the identity 1-morphism on $GF(a) = G(b) = d(c) = d(G(a))$ and its 1-morphism component at 
    $b \in \Adj$ is given by 
    \begin{equation}
        \xymatrix{
            GF(b) = G(a) = c \ar[r]^-{\eta_c} & d^2(c) =d G(b)
        }
    \end{equation}
    Its 2-morphism component at $X \colon a \to b$ in $\Adj$ is given by the identity 2-morphism filling the diagram 
    \begin{equation}
        \begin{aligned}
        \xymatrix{
            GF(a) \ar[r]^-{=} \ar[d]_-{GF(X)} & G(b) \ar[r]^-{=} \ar[d]_-{G(Y) }& d(c) \ar[d]_-{\eta_c^{-1} \circ d(Z)} \ar[r]^-{=} \ar[dr]^-{d(Z)} &dG(a) \ar[dr]^-{dG(X)}\\
            GF(b) \ar[r]_-{=} & G(a) \ar[r]_-{=} & c \ar[r]_-{\eta_c} & d^2(c) \ar[r]_-{=} &dG(b)
        }
        \end{aligned}
    \end{equation}
    Its 2-morphism component at $Y \colon b \to a$ in $\Adj$ is given by 
    \begin{equation}
        \begin{aligned}
        \xymatrix{
            GF(b) \ar[r]^-{=} \ar[d]_-{GF(Y)} & G(a) \ar[r]^-{=} \ar[d]_{G(X)} & c \ar[r]^-{\eta_c} \ar[d]_-{Z} & d^2(c) \ar[r]^-{=} \ar[d]^-{d(\eta_c^{-1} \circ d(Z))} \ar@{=>}[dl]^-{\eta_Z} & dG(b) \ar[d]^-{dG(Y)} \\
            GF(a) \ar[r]_-{=} & G(b) \ar[r]_-{=} & d(c) \ar[r]_-{=} & d(c) \ar[r]_-{=} & dG(a)
        }
        \end{aligned}
    \end{equation}
    One checks that this defines an invertible 2-transformation. Next, we define the invertible modification $\Xi$ fitting into the diagram 
    \begin{equation}
        \begin{aligned}
        \xymatrix{
        ddG \ar[r]^-{\id \circ \omega^{-1}} \ar[dr]_-{\eta \circ \id} & dGF \ar[r]^-{\omega^{-1} \circ \id} & GFF \ar[dl]^-{\id}\\
        & G &
        }
        \end{aligned}
    \end{equation}
    Its 2-morphism component at $a \in \Adj$ is given by the identity 2-morphism on $\eta_c^{-1}$. Its 2-morphism component at $b \in \Adj$ is given 
    by filling the following diagram 
    \begin{equation}
        \begin{aligned}
        \xymatrix{
        ddG(b) \ar[d]^-{=} \ar[r]^-{d(\omega_b^{-1})} & dGF(b) \ar[r]^-{\omega_{F(b)}^{-1}} \ar[d]^-{=} \ar@{=>}[dl]^-{\id} & GFF(b) \ar[d]^-{=} \ar@{=>}[dl]^-{\id} \\
        d^3(c) \ar[dr]_-{\eta_{d(c)}^{-1}} \ar[r]_-{d(\eta_c^{-1})} & d(c) \ar[r]_-{=} \ar@{=>}[d]^-{\tau_c^{-1}} & d(c) \ar[dl]^-{\id} \\
        & G(b) = d(c) &
        }
        \end{aligned}
    \end{equation}
    One checks that the triple $(G,\omega,\Xi)$ defines a 2-contravariant $\Oone$-volutive functor 
    \begin{equation}
        (G,\omega,\Xi) \colon (\Adj,F) \to (\B,d,\eta,\tau).
    \end{equation}
    Conversely, let $(G,\omega,\Xi) \colon \Adj \to \B$ be a 2-contravariant $\Oone$-volutive functor. We define a volutive adjunction 
    whose object is $c := G(a)$, whose 1-morphism $Z \colon c \to d(c)$ is given by the composition 
    \begin{equation}
        Z := \left( \xymatrix{ c = G(a) \ar[r]^-{G(X)} & G(b) = GF(a) \ar[r]^-{\omega_a} & dG(a) = d(c) } \right),
    \end{equation}
    and whose 2-morphism $h \colon \id_c \to \eta_c^{-1} \circ d(Z) \circ Z$ is given by the composition
    \begin{equation}
        \begin{aligned}
        \xymatrix@R=0.5em{
            \id_c = \id_{G(a)} = G(\id_a) \ar[r]^-{G(f)} & G(YX) = G(Y) \circ G(X) \ar[r]^-{=} & G(F(X)) \circ G(X)  \\
            \ar[r]_-{\omega_X \circ \id} & \omega_b^{-1}  dG(X)  \omega_a \omega_a^{-1} Z \ar[r]_-{=} & \omega_b^{-1} dG(X) Z \\
            \ar[r]_-{=} & \omega_b^{-1} d(\omega_a^{-1}\omega_a G(X)) Z \ar[r]_-{=} & \omega_b^{-1} d(\omega_a^{-1})d(\omega_a G(X)) Z\\
            \ar[r]_-{=} & \omega_b^{-1} d(\omega_a^{-1}) d(Z) Z \ar[r]_-{\Xi_c \circ \id} & \eta_c d(Z) Z
        }
        \end{aligned}
    \end{equation}
    One checks that the triple $(c,Z,h)$ defines a volutive adjunction.
\end{proof}

\end{document}
